\BourbakiBackSection{术语索引}

引用编号依次表示章、节、小节（或习题）。

实数的绝对值：IV, 1, 6. 绝对闭空间：I, 9, 习题 18. 绝对收敛无穷乘积：IV, 7, 6. 绝对收敛级数：IV, 7, 6. 可达空间：I, 8, 习题 1. 有理直线的加法群：IV, 1, 2. 实直线的加法群：IV, 1, 3. 拓扑除环的加法一致结构：III, 3, 8. 实直线的加法一致结构：IV, 1, 6. Alexandroff 紧化：I, 9, 8. Alexandroff 半直线：IV, 2, 习题 12. Alexandroff 定理：I, 9, 8. A-极大拓扑：I, 3, 习题 11. 子群的代数补：III, 6, 2. 交错级数：IV, 7, 6. 精确到 ε 的不足（过剩）近似：IV, 8, 1. 拓扑群的任意小子群：III, 2, 习题 30. Archimedes 公理：IV, 2, 1. 相伴双射连续同态：III, 2, 8. 伴随 Hausdorff 群：III, 2, 6. 伴随 Hausdorff 模：III, 6, 6. 伴随 Hausdorff 环：III, 6, 4. 结合性公式：III, 5, 3. 可和族之和的结合性：III, 5, 3. 自同构，局部：III, 1, 3. 拓扑群的自同构：III, 1, 3. Archimedes 公理：IV, 2, 1. Borel-Lebesgue 公理：I, 9, 1. Hausdorff 公理：I, 8, 1.

\begin{Shaded}
\begin{Highlighting}[]
\NormalTok{滤子基：I, 6, 3.}
\NormalTok{实数展开式的底（数）：IV, 8, 5.}
\NormalTok{拓扑基：I, 1, 3.}
\NormalTok{实数展开式的基序列：IV, 8, 2.}
\NormalTok{双连续映射：I, 2, 1.}
\NormalTok{Bolzano\textquotesingle{}s 定理：IV, 6, 1.}
\NormalTok{Borel{-}Lebesgue 公理：I, 9, 1.}
\NormalTok{Borel{-}Lebesgue 定理：IV, 2, 2.}
\NormalTok{上有界、下有界（函数）：IV, 5, 1.}
\NormalTok{上有界、下有界（集合）：IV, 2, 3.}
\NormalTok{有界函数：IV, 5, 1.}
\NormalTok{拓扑环的左（右）有界子集：III, 6, 习题 12.}
\NormalTok{有界集（一致空间中的）：II, 4, 习题 7.}
\NormalTok{函数的下确界与上确界：IV, 5, 9.}

\NormalTok{由自由作用的群所定义的等价关系的图的典范映射：III, 4, 3.}
\NormalTok{Cantor\textquotesingle{}s 定理：IV, 8, 6.}
\NormalTok{Cantor\textquotesingle{}s 三分集：IV, 2, 5.}
\NormalTok{Cauchy 滤子：II, 3, 1.}
\NormalTok{Cauchy 滤子，极小：II, 3, 2.}
\NormalTok{Cauchy 序列：II, 3, 1.}
\NormalTok{Cauchy\textquotesingle{}s 并项判别法：IV, 7, 习题 3.}
\NormalTok{Cauchy\textquotesingle{}s 判别准则：II, 3, 3.}
\NormalTok{级数的 Cauchy\textquotesingle{}s 判别准则：III, 5, 46.}
\NormalTok{可和族的 Cauchy\textquotesingle{}s 判别准则：III, 5, 2.}
\NormalTok{链（V{-}）：II, 4, 4.}
\NormalTok{求和次序的交换：III, 5, 3.}
\NormalTok{V{-} 邻近：II, 1, 1.}
\NormalTok{闭覆盖：I, 1, 4.}
\NormalTok{闭等价关系：I, 5, 2.}
\NormalTok{闭映射：I, 5, 1.}
\NormalTok{闭集：I, 1, 4.}
\NormalTok{集合的闭包：I, 1, 6.}
\NormalTok{滤子基的聚点：I, 7, 2.}
\NormalTok{函数相对于子集在一点的聚点：I, 7, 5.}
\NormalTok{函数关于有向集的聚点：I, 7, 3.}
\NormalTok{函数关于滤子的聚点：I, 7, 3.}
\NormalTok{函数关于序列的聚点：I, 7, 3.}
\NormalTok{函数的芽的聚点：I, 7, 3.}
\NormalTok{较粗滤子：I, 6, 2.}
\NormalTok{较粗拓扑：I, 2, 2.}
\NormalTok{较粗一致结构：II, 2, 2.}
\NormalTok{最粗拓扑：I, 2, 2.}
\end{Highlighting}
\end{Shaded}

交换收敛级数：III, 5, 7. 紧集：I, 9, 3. 紧空间：I, 9, 3. 紧化，Alexandroff 或单点：I, 9, 8. 可比较的滤子：I, 6, 2. 可比较的拓扑：I, 2, 2. 可比较的一致结构：II, 2, 2. 级数的比较原理：IV, 7, 1. 相容（除环结构与拓扑）：III, 6, 7. 相容（等价关系与算子群）：III, 2, 4. 相容（群结构与拓扑）：III, 1, 1. 相容（带算子空间的映射与算子群的同态）：III, 2, 4. 相容（有序群结构与拓扑）：IV, 1, 习题 1. 相容（环结构与拓扑）：III, 6, 3. 相容（带算子的群结构与拓扑）：III, 6, 1. 与拓扑相容：II, 4, 1. 补，代数：III, 6, 2. 补，拓扑：III, 6, 2. 完备群：III, 3, 3. 完备环：III, 6, 5. 完备一致空间：II, 3, 3. 完全 Hausdorff（空间、拓扑）：I, 9, 习题 20. 两点间完全不可约：IV, 2, 习题 16. Hausdorff 拓扑除环、域的完备化：III, 6, 8. Hausdorff 拓扑群的完备化：III, 3, 4. Hausdorff 拓扑模的完备化：III, 6, 6. Hausdorff 拓扑环的完备化：III, 6, 5. Hausdorff 一致空间的完备化：II, 3, 7. 分支，单位元：III, 2, 2. 点的分支：I, 11, 5. 子集的分支：I, 11, 5. 凝聚点：I, 9, 习题 16. 并项判别法，Cauchy 的：IV, 7, 习题 3. 连通分支：I, 11, 5. 连通集：I, 11, 1. 连通空间：I, 11, 1. 组成部分，素：II, 4, 习题 19. 邻接区间：IV, 2, 5. 借连续性的延拓：I, 8, 5. 与拓扑群的连续同态相伴的双射连续同态：III, 2, 8. 连续（映射、函数）：I, 2, 1. 连续截面：I, 3, 5.

关于子空间连续：I, 3, 2. 连续地，群的作用：III, 2, 4. 连续统的势：IV, 8, 6. 收敛，绝对：IV, 7, 6. 收敛，交换：III, 5, 7. 收敛滤子：I, 7, 1. 收敛滤子基：I, 7, 1. 收敛无穷乘积：III, 5, 7. 实数的收敛无穷乘积：IV, 7, 6. 收敛序列：I, 7, 3. 收敛级数：III, 5, 6. 实数的收敛级数：IV, 7, 6. 余预滤子：I, 6, 习题 17. 对应，逆紧：I, 10, 习题 10. 覆盖，闭：I, 1, 4. 覆盖，开：I, 1, 1. 判别准则，Cauchy 的：II, 3, 3；III, 5, 2 与 III, 5, 6. 立方根：IV, 3, 3.

实数的十进展开：IV, 8, 5. 稠密集：I, 1, 6. 直和，拓扑：III, 6, 2. Dirichlet 函数：IV, 6, 2. 离散除环：III, 6, 7. 离散域：III, 6, 7. 离散群：III, 1, 1. 离散环：III, 6, 3. 离散（拓扑空间、拓扑）：I, 1, 1. 离散（一致空间、一致结构）：II, 1, 1. 除环，拓扑：III, 6, 7. 实数的二进展开：IV, 8, 5.

初等滤子：I, 6, 8. 与序列相伴的初等滤子：I, 6, 8. 初等集：I, 4, 1. 局部紧空间的端：I, 11, 习题 19. 一致邻域：II, 1, 1. 一致邻域，对称：II, 1, 1. 一致邻域基本系：II, 1, 1. 包络，下（上）：IV, 5, 5. 等价关系，闭：I, 5, 2. 由算子群定义的等价关系：III, 2, 4. 等价关系，Hausdorff：I, 8, 3. 等价关系，开：I, 5, 2.

展开，十进：IV, 8, 5.

展开，二进：IV, 8, 5.

展开，非正常：IV, 8, 3.

实数关于基序列的展开：IV, 8, 2.

展开，有限（终止）：IV, 8, 3.

以 a 为底的展开：IV, 8, 5.

展开，三进：IV, 8, 5.

扩充实直线：IV, 4, 2.

映射借连续性的延拓：I, 8, 5.

恒等式延拓原理：I, 8, 1.

不等式延拓原理：IV, 5, 2.

集合的外部：I, 1, 6.

外部点：I, 1, 6.

外半直积：III, 2, 10.

外拓扑半直积：III, 2, 10.

极端不连通空间：I, 11, 习题 21.

因子，通（无穷乘积的）：III, 5, 7.

族，局部有限：I, 1, 5.

族，可乘：III, 5, 1.

族，可和：III, 5, 1.

域，离散：III, 6, 7.

实数域：IV, 3, 1.

域，p 进：III, 6, 习题 23.

域，拓扑：III, 6, 7.

滤子：I, 6, 1.

滤子基：I, 6, 3.

滤子，Cauchy：II, 3, 1.

较粗于另一滤子的滤子：I, 6, 2.

与另一滤子可比较的滤子：I, 6, 2.

滤子，收敛：I, 7, 1.

滤子，初等：I, 6, 8.

滤子，初等，与序列相伴的：I, 6, 8.

较细于另一滤子的滤子：I, 6, 2.

滤子，Fréchet：I, 6, 1.

由子集族生成的滤子：I, 6, 2.

滤子，诱导：I, 6, 5.

滤子，交：I, 6, 2.

滤子，极小 Cauchy：II, 3, 2.

滤子，邻域：I, 6, 1.

滤子，截面：I, 6, 3.

严格粗于另一滤子的滤子：I, 6, 2.

严格细于另一滤子的滤子：I, 6, 2.

滤子子基：I, 6, 2.

滤子化集：I, 6, 1.

终拓扑：I, 2, 4.

较细滤子：I, 6, 2.

较细拓扑：I, 2, 2.

较细一致结构：II, 2, 2.

有限开覆盖的一致结构：II, 4, 1.

有限部分和：III, 5, 1.

有限分拆的一致结构：II, 2, 2.

有限实数：IV, 4, 2.

有限实值函数：IV, 5, 1.

自由地，群的作用：III, 4, 3.

Fréchet 滤子：I, 6, 1.

集合的边界：I, 1, 6.

边界点：I, 1, 6.

全子集：I, 11, 习题 14.

函数，连续：I, 2, 1.

函数，实值（或实）：IV, 5, 1.

函数，一致连续：II, 2, 1.

一致邻域基本系：II, 1, 1.

基本邻域系：I, 1, 3.

Gelfand 环：III, 6, 习题 11.

无穷乘积的通因子：III, 5, 7.

级数的通项：VI, 5, 6.

等比数列（几何级数）：IV, 7, 1.

映射在一点的芽：I, 6, 10.

映射关于滤子的芽：I, 6, 9.

子集在一点的芽：I, 6, 10.

子集关于滤子的芽：I, 6, 9.

实值函数的下确界：IV, 5, 4.

拓扑族的下确界：I, 2, 4.

一致结构族的下确界：II, 2, 5.

群，完备：III, 3, 3.

群，离散：III, 1, 1.

群，Hausdorff，与拓扑群相伴的：III, 2, 6.

自由作用于集合上的群：III, 4, 3.

平凡作用于集合上的群：III, 2, 4.

群，仿拓扑：III, 1, 习题 4.

群，积拓扑：III, 2, 9.

群，拟拓扑：III, 2, 习题 5.

群，商拓扑：III, 2, 6.

群，半拓扑：III, 1, 习题 2.

群，拓扑：III, 1, 1.

群，拓扑，连续作用于拓扑空间上的：III, 2, 4.

群，拓扑，逆紧作用于拓扑空间上的：III, 4, 1. 群，拓扑，带算子的：III, 6, 1.

半直线，Alexandroff：IV, 2, 习题 12. Hamel 基：IV, 6, 习题 1. 拓扑群的 Hausdorff 完备化：III, 3, 4. 拓扑模的 Hausdorff 完备化：III, 6, 6. 拓扑环的 Hausdorff 完备化：III, 6, 5. 一致空间的 Hausdorff 完备化：II, 3, 7. 等价关系，Hausdorff：I, 8, 3. 与拓扑群相伴的 Hausdorff 群：III, 2, 6. 与拓扑模相伴的 Hausdorff 模：III, 6, 6. 与拓扑环相伴的 Hausdorff 环：III, 6, 4. Hausdorff 空间：I, 8, 1. Hausdorff 拓扑：I, 8, 1. 与一致空间相伴的 Hausdorff 一致空间：II, 3, 8. Hausdorff 一致结构：II, 1, 2. Hausdorff 公理：I, 8, 1. 同胚的拓扑空间：I, 1, 1. 同胚：I, 1, 1. 同胚，局部：I, 11, 习题 25. 齐性空间：III, 2, 5.

借粘合得到的空间：I, 3, 4. 恒等式延拓原理：I, 8, 1. 拓扑群的单位元连通分支：III, 2, 2. 拓扑的逆像：I, 2, 3. 一致结构的逆像：II, 2, 4. 实数的非正常展开：IV, 8, 3. 由一致结构诱导的（拓扑）：II, 1, 2. 诱导滤子：I, 6, 5. 诱导拓扑：I, 2, 3. 诱导一致结构：II, 2, 4. 无穷乘积，绝对收敛：IV, 7, 6. 无穷乘积，收敛：III, 5, 7. 由序列 ( \(x_n\) ) 定义的无穷乘积：III, 5, 7. 实数的无穷乘积，收敛：IV, 7, 6. 通因子为 \(x_n\) 的无穷乘积：III, 5, 7. 初始拓扑：I, 2, 3. 初始一致结构：II, 2, 3. 拓扑群的内自同构：III, 1, 3. 整数，p 进：III, 6, 习题 23. 实数的整数部分：IV, 8, 2. 集合的内部：I, 1, 6.

内部点：I, 1, 6.

交滤子：I, 6, 2.

交拓扑：I, 1, 6.

拓扑的逆像：I, 2, 3.

一致结构的逆像：II, 2, 4.

群（环）的逆系统的逆极限：III, 7, 1.

赋有内（外）合成法则的集合的逆系统的逆极限：III, 7, 1.

拓扑群（模、环）的逆系统的逆极限：III, 7, 2.

拓扑空间的逆极限：I, 4, 4.

拓扑的逆极限：I, 4, 4.

一致空间的逆极限：II, 2, 7.

一致结构的逆极限：II, 2, 7.

群（环）的逆系统：III, 7, 1.

拓扑群（环）的逆系统：III, 7, 2.

拓扑空间的逆系统：I, 4, 4.

拓扑的逆系统：I, 4, 4.

一致空间的逆系统：II, 2, 7.

一致结构的逆系统：II, 2, 7.

无理数：IV, 1, 2.

两点间不可约：II, 4, 习题 19.

同密空间：I, 2, 习题 11.

孤立点：I, 1, 6.

同构的一致空间：II, 1, 1.

同构，局部：III, 1, 3.

拓扑群到拓扑群上的同构：III, 1, 3.

一致空间到另一个一致空间上的同构：II, 1, 1.

Kolmogoroff 空间：I, 1, 习题 2.

实值函数的上确界：IV, 5, 4.

拓扑族的上确界：I, 2, 3.

一致结构族的上确界：II, 2, 5.

左附着点：IV, 2, 习题 1.

左有界：III, 6, 习题 12.

线性映射的左逆：III, 6, 2.

有序集上的左拓扑：I, 1, 习题 2.

拓扑群的左一致结构：III, 3, 1.

有界区间的长度：IV, 1, 5.

极限，逆：III, 7, 1 与 2.

极限，下（上）：IV, 5, 6.

左（右）极限：IV, 5, 3.

滤子的极限（点）：I, 7, 1.

滤子基的极限（点）：I, 7, 1.

函数在一点的极限（点）：I, 7, 5.

函数相对于子集的极限（点）：I, 7, 5.

函数关于有向集的极限（点）：I, 7, 3.

函数关于滤子的极限（点）：I, 7, 3.

函数的芽的极限（点）：I, 7, 3.

序列的极限（点）：I, 7, 3.

Lindelöf 空间：I, 9, 习题 14.

直线，扩充实：IV, 4, 2.

直线，有理：I, 1, 2；IV, 1, 2.

直线，实：IV, 1, 3.

拓扑群的局部自同构：III, 1, 3.

拓扑群的局部直积：III, 2, 习题 26.

局部同胚：I, 11, 习题 25.

拓扑群与另一拓扑群的局部同构：III, 1, 3.

拓扑群族关于开正规子群族的局部积：III, 2, 习题 26.

局部有界拓扑环：III, 6, 习题 12.

局部闭子空间：I, 3, 2.

局部紧空间：I, 9, 7.

局部连通空间：I, 11, 6.

局部有限族：I, 1, 5.

局部同构的拓扑群：III, 1, 3.

局部准紧 Hausdorff 拓扑群：III, 3, 习题 8.

局部拟紧空间：I, 9, 习题 29.

局部逆有界拓扑：III, 6, 习题 22.

实值函数族的下包络：IV, 5, 5.

实值函数关于滤子的下极限：IV, 5, 6.

下半连续实值函数：IV, 6, 2.

实值函数的下半连续正则化：IV, 6, 2.

映射，双连续：I, 2, 1.

映射，闭：I, 5, 1.

映射，连续：I, 2, 1.

映射，开：I, 5, 1.

映射，逆紧：I, 10, 1.

映射，一致连续：II, 2, 1.

极大，相对：IV, 6, 2.

极大，严格相对：IV, 5, 习题 10.

极小 Cauchy 滤子：II, 3, 2.

极小 Hausdorff 空间：I, 9, 习题 19.

极小，相对：IV, 6, 2.

Mittag-Leffler 定理：II, 3, 5.

\begin{Shaded}
\begin{Highlighting}[]
\NormalTok{模，拓扑：III, 6, 6.}
\NormalTok{带算子空间的态射：III, 2, 4.}
\NormalTok{拓扑群的态射：III, 2, 8.}
\NormalTok{态射，严格：III, 2, 8.}
\NormalTok{单调极限定理：IV, 5, 2.}
\NormalTok{乘法群中元素的可乘族：III, 5, 1.}
\NormalTok{拓扑除环的乘法一致结构：III, 6, 8.}

\NormalTok{邻域滤子：I, 6, 1.}
\NormalTok{点的邻域：I, 1, 2.}
\NormalTok{集合的邻域：I, 1, 2.}
\NormalTok{邻域，对称：III, 1, 2.}
\NormalTok{邻域（V{-}）：II, 1, 2.}
\NormalTok{邻域系，基本：I, 1, 3.}
\NormalTok{幂零，拓扑上：III, 6, 习题 9.}
\NormalTok{n 次根：IV, 3, 3.}
\NormalTok{数，有限实：IV, 5, 2.}
\NormalTok{数，无理：IV, 1, 3.}
\NormalTok{数，p{-}进：III, 6, 习题 23.}
\NormalTok{数，实：IV, 1, 3 及（滥用术语）IV, 4, 2.}

\NormalTok{单点{-}紧化：I, 9, 8.}
\NormalTok{开覆盖：I, 1, 1.}
\NormalTok{开等价关系：I, 5, 2.}
\NormalTok{开映射：I, 5, 1.}
\NormalTok{开集：I, 1, 1.}
\NormalTok{拓扑群的反群：III, 1, 1.}
\NormalTok{带算子空间关于其算子群的轨道空间：III, 2, 4.}
\NormalTok{实值{-}函数的振荡：IV, 6, 习题 6.}

\NormalTok{p{-}进域、整数、数、拓扑、赋值：III, 6, 习题 23.}
\NormalTok{p{-}进螺线管：III, 7, 习题 6.}
\NormalTok{p{-}进一致结构：II, 1, 1.}
\NormalTok{仿紧空间：I, 9, 10.}
\NormalTok{仿拓扑群：III, 1, 习题 4.}
\NormalTok{部分和：III, 5, 3.}
\NormalTok{集合的粘合：I, 2, 5.}
\NormalTok{拓扑空间的粘合：I, 2, 5.}
\NormalTok{完备集：I, 1, 6.}
\NormalTok{实数的循环展开：IV, 8, 习题 3.}
\NormalTok{可交换的（两个群在同一空间上的作用）：III, 2, 4.}
\NormalTok{Poincaré{-}Volterra 定理：I, 11, 7.}
\NormalTok{点：I, 1, 1.}
\NormalTok{点，聚：I, 7, 2.}
\end{Highlighting}
\end{Shaded}

点，凝聚：I, 9, 习题 16. 点，外部：I, 1, 6. 点，边界：I, 1, 6. 点，内部：I, 1, 6. 点，孤立：I, 1, 6. 点，极限：I, 7, 1. 点，奇点：II, 4, 习题 19. 点，V-邻近：II, 1, 1. 连续统的势：IV, 8, 6. 准紧集：II, 4, 2. 准紧空间：II, 4, 2. 预滤子：I, 6, 习题 17. 素组成部分：II, 4, 习题 19. 素组成部分的空间：II, 4, 习题 19. 素预滤子：I, 6, 习题 17. 本原集：I, 7, 习题 8. 级数的比较原理：IV, 7, 1. 恒等式延拓原理：I, 8, 1. 不等式延拓原理：IV, 5, 2. 积，外半直：III, 2, 10. 积滤子：I, 6, 7. 积，局部（直）：III, 2, 习题 26. 可乘族的积：III, 5, 1. 拓扑群的积：III, 2, 9. 带算子拓扑群的积：III, 6, 1. 拓扑环的积：III, 6, 4. 积，半直：III, 2, 10. 积，拓扑半直：III, 2, 10. 积拓扑空间：I, 4, 1. 积拓扑：I, 4, 1. 积一致空间：II, 2, 6. 积一致结构：II, 2, 6. 对应，逆紧：I, 10, 习题 10. 映射，逆紧：I, 10, 1. 逆紧地，群的作用：III, 4, 1.

拟紧集：I, 9, 3. 拟紧空间：I, 9, 1. 拟极大拓扑：I, 2, 习题 6. 拟环：III, 6, 习题 20. 拟拓扑群：III, 2, 习题 5. 商空间：I, 3, 4. 带算子空间关于其算子群的商空间：III, 2, 4. 商拓扑群：III, 2, 6.

带算子的商拓扑群：III, 6, 1.

商拓扑环：III, 6, 4.

商拓扑：I, 3, 4.

群的拓扑关于子群的商拓扑：III, 2, 5.

带算子空间的拓扑关于其算子群的商拓扑：III, 2, 4.

交换拓扑群的根：III, 2, 习题 28.

根群：III, 2, 习题 28.

有理直线：I, 1, 2；IV, 1, 2.

n 维有理数空间：I, 4, 1.

实函数：IV, 5, 1.

实直线：IV, 1, 3.

实直线，扩充实直线：IV, 4, 2.

实数：IV, 1, 3 及（滥用术语）IV, 4, 2.

实值函数：IV, 5, 1.

实值函数，上（下）有界的：IV, 5, 4.

实值函数，有限：IV, 5, 1.

实值函数，下（上）半连续：IV, 6, 2.

正则开集：I, 8, 习题 20.

正则空间：I, 8, 4.

正则拓扑：I, 8, 4.

正则化，下（上）半连续：IV, 6, 2.

相对极大（极小）：IV, 6, 2.

相对紧：I, 9, 3.

相对拟紧：I, 9, 3.

剩余闭子群：III, 2, 习题 28.

级数的余项：III, 5, 6.

级数的受限结合性：III, 5, 7.

逆有界集：III, 6, 习题 22.

右有界：III, 6, 习题 12.

线性映射的右逆：III, 6, 2.

右拓扑（有序集上的）：I, 1, 习题 2.

拓扑群的右一致结构：III, 3, 1.

环，完备：III, 6, 5.

环，离散：III, 6, 3.

环，Gelfand：III, 6, 习题 11.

环，Hausdorff，与拓扑环相伴的：III, 6, 4.

环，局部有界拓扑：III, 6, 习题 12.

环，积拓扑：III, 6, 4.

环，商拓扑：III, 6, 4.

环，拓扑：III, 6, 3.

环，拓扑，Hausdorff 完备化：III, 6, 5.

根（平方根、立方根、n 次根）：IV, 3, 3.

σ-紧：I, 9, 9.

子集关于算子群的饱和集：III, 2, 4.

截面，连续：I, 3, 5.

截面滤子：I, 6, 3.

半连续，下（上）：IV, 6, 2.

子群的半直积：III, 2, 10.

半拓扑群：III, 1, 习题 2.

半正则（空间、拓扑）：I, 8, 习题 20.

序列，Cauchy：II, 3, 1.

收敛序列：I, 7, 3.

级数，绝对收敛：IV, 7, 6.

级数，交错：IV, 7, 6.

级数，交换收敛：III, 5, 7.

级数，收敛：III, 5, 6 与 IV, 7, 6.

由序列 \((x_{n})\) 定义的级数：III, 5, 6.

级数，其通项为 \(x_{n} :\) III, 5, 6.

集，有界：II, 4, 习题 7.

集，上（下）有界：IV, 2, 3.

集，Cantor 三分集：IV, 2, 5.

集，闭：II, 1, 4.

集，紧：I, 9, 3.

集，连通：I, 11, 1.

集，稠密：I, 1, 6.

集，初等：I, 4, 1.

集，滤子化：I, 6, 1.

集，局部闭：I, 3, 3.

集，开：I, 1, 1.

集，完备：I, 1, 6.

集，准紧：II, 4, 2.

集，拟紧：I, 9, 3.

集，正则开：I, 8, 习题 20.

集，相对紧：I, 9, 3.

集，相对拟紧：I, 9, 3.

集，全不连通：I, 11, 5.

拓扑空间的底集：I, 1, 1.

实数的符号：IV, 3, 2.

奇点：II, 4, 习题 19.

小（V-）：II, 3, 1.

螺线管，p 进：III, 7, 习题 6.

可解空间：I, 2, 习题 11.

空间，绝对闭：I, 9, 习题 18.

空间，可达：I, 8, 习题 i.

空间，紧：I, 9, 1.

空间，完备：II, 3, 3.

空间，完全 Hausdorff：I, 9, 习题 20.

空间，连通：I, 11, 1.

空间，极端不连通：I, 11, 习题 21.

空间，Hausdorff：I, 8, 1.

空间，Hausdorff，与一致空间相伴的：II, 3, 8.

空间，齐性：III, 2, 5.

空间，同密：I, 2, 习题 11.

空间，Kolmogoroff：I, 1, 习题 2.

空间，Lindelöf：I, 9, 习题 14.

空间，局部紧：I, 9, 7.

空间，局部紧 \(\sigma\) -紧：I, 9, 9.

空间，局部连通：I, 11, 6.

空间，局部拟紧：I, 9, 习题 29.

空间，极小 Hausdorff：I, 9, 习题 19.

空间， \(n\) 维有理：I, 4, 1.

连续作用于拓扑空间上的群的轨道空间：III, 2, 4.

素组成部分的空间：II, 4, 习题 19.

空间，仿紧：I, 9, 10.

空间，准紧：II, 4, 2.

空间，积拓扑：I, 4, 1.

空间，积一致：II, 2, 6.

空间，拟紧：I, 9, 1.

空间，商：I, 3, 4；III, 2, 4.

空间，n 维有理：I, 4, 1.

空间，正则：I, 8, 4.

空间，半正则：I, 8, 习题 20.

空间，可解：I, 2, 习题 11.

空间，次极大：I, 8, 习题 22.

空间，和：I, 2, 4.

空间，拓扑：I, 1, 1.

空间，拓扑齐性：III, 2, 5.

空间，拓扑，一致空间的底拓扑空间：II, 1, 2.

空间，全不连通：I, 11, 5.

空间，超滤子：I, 9, 习题 26.

空间，超正则：I, 8, 习题 25.

空间，一致：II, 1, 1.

空间，可一致化：II, 4,1.

平方根：IV, 3, 3.

稳定子群：III, 2, 4.

严格态射：III, 2, 8.

严格相对极大：IV, 5, 习题 10.

严格较粗滤子：I, 6, 2.

严格较粗拓扑：I, 2, 2.

严格较粗一致结构：II, 2, 2.

严格较细滤子：I, 6, 2.

严格较细拓扑：I, 2, 2.

严格较细一致结构：II, 2, 2.

滤子的子基：I, 6, 2.

拓扑的子基：I, 2, 3.

次极大（空间、拓扑）：I, 8, 习题 22.

子空间（拓扑空间的）：I, 3, 1.

子空间，局部闭：I, 3, 3.

一致空间的子空间：II, 2, 4.

和，有限部分：III, 5, 1.

交换拓扑群中点族的和：III, 5, 1.

级数的和：III, 5, 6.

和，部分：III, 5, 3.

和（拓扑空间、拓扑）：I, 2, 4.

可和族：III, 5, 1.

对称一致邻域：II, 1, 1.

对称邻域：III, 1, 2.

对称映射：III, 1, 1.

有限（终止）展开（实数的）：IV, 8, 3.

定理，Alexandroff 的：I, 9, 8.

定理，Mittag-Leffler 的：II, 3, 5.

定理，Poincaré-Volterra 的：I, 11, 7.

定理，Tychonoff 的：I, 9, 5.

Bolzano 定理：IV, 6, 1.

Borel-Lebesgue 定理：IV, 2, 2.

Cantor 定理：IV, 8, 6.

单调极限定理：IV, 5, 2.

Weierstrass 定理：IV, 6, 1.

稳定子群的拓扑补：III, 6, 2.

稳定子群的拓扑直和：III, 6, 2.

拓扑除环：III, 6, 7.

拓扑域：III, 6, 7.

拓扑群：III, 1, 1.

带算子拓扑群：III, 6, 1.

拓扑齐性空间：III, 2, 5.

拓扑模：III, 6, 6.

拓扑环：III, 6, 3.

拓扑半直积：III, 2, 10.

拓扑空间：I, 1, 1.

拓扑结构：见拓扑.

拓扑幂零（元素、理想）：III, 6, 习题 9.

拓扑：I, 1, 1. 拓扑，A-极大：I, 3, 习题 11. 与滤子相伴的拓扑：I, 6, 5. 较另一拓扑粗的拓扑：I, 2, 2. 拓扑，最粗：I, 2, 2. 与另一拓扑可比较的拓扑：I, 2, 2. 拓扑，完全 Hausdorff：I, 9, 习题 20. 拓扑，离散：I, 1, 1. 拓扑，终：I, 2, 4. 较另一拓扑细的拓扑：I, 2, 2. 由子集族生成的拓扑：I, 2, 3. 拓扑，Hausdorff：I, 8, 1. 拓扑，诱导：I, 2, 3. 由一致结构诱导的拓扑：II, 1, 2. 拓扑，初始：I, 2, 3. 拓扑，交：I, 1, 6. 拓扑，左（有序集上的）：I, 1, 习题 2. 拓扑，局部有界：III, 6, 习题 12. 拓扑，局部逆有界：III, 6, 习题 22. 拓扑，p 进：III, 6, 习题 23. 拓扑，积：I, 2, 3. 拓扑，拟极大：I, 2, 习题 6. 拓扑，商：I, 2, 4. 拓扑，正则：I, 8, 4. 拓扑，右（有序集上的）：I, 1, 习题 2. 拓扑，半正则，与拓扑相伴的：I, 8, 习题 20. 严格较另一拓扑粗的拓扑：I, 2, 2. 严格较另一拓扑细的拓扑：I, 2, 2. 拓扑，次极大：I, 8, 习题 22. 拓扑，和：I, 2, 4. 拓扑，超正则：I, 8, 习题 25. 拓扑，可一致化：II, 4, 1. 全附着：I, 9, 习题 17. 全不连通（集、空间）：I, 11, 5. 实数的三进展开：IV, 8, 5. 三分集，Cantor 的：IV, 2, 6. 三角不等式：IV, 1, 1. 平凡超滤子：I, 6, 4. 平凡地，群作用：III, 2, 4. 拓扑群上的双侧一致结构：III, 2, 习题 6. Tychonoff 定理：I, 9, 5.

超滤子：I, 6, 4

超滤子空间：I, 9, 习题 26.

超滤子，平凡：I, 6, 4.

超正则（空间、拓扑）：I, 8, 习题 25.

拓扑空间的底集：I, 1, 1.

一致空间的底拓扑空间：II, 1, 2.

一致结构：见一致结构.

一致结构，乘法（拓扑除环的）：III, 6, 8.

一致结构：II, 1, 1.

一致结构，加法（拓扑除环的）：III, 6, 8.

一致结构，加法（实直线的）：IV, 1, 6.

较另一致结构粗的一致结构：II, 2, 2.

与另一致结构可比较的一致结构：II, 2, 2.

与拓扑相容的一致结构：II, 4, 1.

一致结构，离散：II, 1, 1.

较另一致结构细的一致结构：II, 2, 2.

一致结构，Hausdorff：II, 1, 2.

一致结构，诱导：II, 2, 4.

一致结构，左（右）（拓扑群的）：III, 3, 1.

有限开覆盖的一致结构：II, 4, 1.

有限分拆的一致结构：II, 2, 2.

一致结构，p 进：II, 1, 1.

一致结构，积：II, 2, 6.

严格较另一致结构粗的一致结构：II, 2, 2.

严格较另一致结构细的一致结构：II, 2, 2.

一致结构，双侧：III, 3, 习题 6.

可一致化（拓扑空间、拓扑）：II, 4, 1.

一致上有界（下有界）（实值函数族）：IV, 5, 5.

一致连续映射：II, 2, 1.

实值函数族的上包络：IV, 5, 5.

实值函数关于滤子的上极限：IV, 5, 5.

上半连续实值函数：IV, 6, 2.

实值函数的上半连续正则化：IV, 6, 2.

赋值，p 进：III, 6, 习题 23.

值，绝对：IV, 1, 6.

映射在一点的芽的值：I, 6, 10.

V-链：II, 4, 4.

V-邻近（点）：II, 1, 1.

集合的 V-邻域：II, 1, 2.

V-小：II, 3, 1.

Weierstrass 定理：IV, 6, 1.
